\documentclass[10pt,twocolumn,letterpaper]{article}

\usepackage{graphicx}
\usepackage{amsmath}
\usepackage{amssymb}
\usepackage{booktabs}
\usepackage[accsupp]{axessibility} % Improves accessibility of PDF
\usepackage{caption}
\usepackage{subcaption}
\usepackage{hyperref}

% Define a shorthand for et al.
\newcommand{\etal}{\textit{et al.}}

% --- TITLE AND AUTHOR INFORMATION ---
\title{\textbf{Architectural Amalgamation in Single Image Super-Resolution: A Comparative Study of Degradation Modeling and Self-Supervised Features}}

\author{
    [Your Name(s)]\\
    [Your Affiliation(s)]\\
    {\tt\small [your.email@example.com]}
}

\date{} % No date to be shown

\begin{document}

\maketitle

% =================================================================================================
% 1. ABSTRACT
% =================================================================================================
\begin{abstract}
\noindent Single Image Super-Resolution (SISR) aims to restore high-resolution details from a single low-resolution input, a fundamentally ill-posed problem. Modern approaches have focused on sophisticated network architectures and training strategies to bridge the performance gap, particularly for real-world images with complex, unknown degradations. This paper presents a comparative study of four distinct architectural paradigms, systematically exploring the amalgamation of three key components: (1) realistic, on-the-fly degradation modeling, (2) advanced Transformer-based backbones, and (3) novel loss functions, including a self-supervised perceptual loss (DINO) and a frequency-domain (FFT) loss. Our four proposed models—\texttt{Degraded}, \texttt{Proposed-Merge}, \texttt{DINO-Ranger}, and \texttt{DINO w/ Degraded}—are trained on the DF2K dataset and evaluated on four standard benchmarks (Set5, Set14, BSD100, Urban100) for x4 super-resolution. 

Experimental results reveal significant insights into the synergies and trade-offs of these components. The \texttt{DINO-Ranger} model, which combines a Hybrid Swin Transformer with DINO and FFT losses, achieves state-of-the-art performance, yielding an average PSNR of 27.51 dB across all benchmarks. Notably, it improves PSNR by up to 3.09 dB over the degradation-focused baseline on the structured Urban100 dataset. Our analysis further shows that while explicit degradation modeling enhances visual naturalness (NIQE metric), it can conflict with powerful perceptual losses, slightly capping fidelity. Conversely, self-supervised features from DINO provide a substantial boost in both quantitative fidelity (PSNR) and perceptual quality (LPIPS), demonstrating their efficacy as a replacement for traditional VGG-based losses. This work provides a clear, evidence-based guide for designing robust and high-performing SISR architectures.
\end{abstract}

% =================================================================================================
% 2. INTRODUCTION
% =================================================================================================
\section{Introduction}
\label{sec:intro}

Single Image Super-Resolution (SISR) is a classical computer vision task with wide-ranging applications, from enhancing medical imagery and satellite surveillance to improving consumer photography. The core objective is to reconstruct a high-resolution (HR) image from its low-resolution (LR) counterpart. This process is inherently ill-posed, as a single LR image can be mapped to multiple potential HR images.

Early methods relied on simple interpolation techniques like bicubic or Lanczos resizing, which often produce blurry and over-smoothed results. The advent of deep learning, particularly Convolutional Neural Networks (CNNs), revolutionized the field. Models like SRCNN \cite{dong2015image} and EDSR \cite{lim2017enhanced} demonstrated the power of deep networks in learning the complex mapping from LR to HR space. More recently, Vision Transformers (ViTs) \cite{dosovitskiy2020image} and their variants, such as SwinIR \cite{liang2021swinir}, have pushed the boundaries further by leveraging self-attention mechanisms to capture long-range dependencies, proving highly effective for image restoration.

Despite these architectural advancements, a significant challenge remains: the "domain gap" between synthetic training data and real-world images. Most models are trained on LR images generated via simple bicubic downsampling. Real-world images, however, suffer from a cocktail of degradations, including sensor noise, motion blur, compression artifacts, and chromatic aberration. Models trained without exposure to these effects often fail to generalize well.

This observation has motivated two parallel research directions: (1) incorporating realistic degradation models into the training pipeline, as popularized by Real-ESRGAN \cite{wang2021real}, and (2) employing more powerful, perceptually-aligned loss functions that guide the model to produce visually plausible results. Traditional perceptual losses, based on pre-trained VGG networks \cite{simonyan2014very}, have been effective but are limited by their original training objective (image classification). The rise of self-supervised learning offers a compelling alternative. Models like DINO \cite{caron2021emerging}, trained without human labels, learn rich semantic and textural features that are potentially better aligned with human perception.

In this paper, we conduct a systematic comparative study to understand the interplay between these modern components. We propose and evaluate four model variants, each representing a unique amalgamation of architectural choices, degradation modeling, and loss functions. Our contributions are as follows:
\begin{itemize}
    \item We implement and compare four SISR models, exploring combinations of a novel degradation pipeline, two distinct Transformer-based generators, and an advanced loss suite featuring DINO and FFT losses.
    \item We provide a detailed quantitative and qualitative analysis on four standard benchmarks for x4 super-resolution, evaluating fidelity (PSNR, SSIM), perceptual quality (LPIPS), and no-reference naturalness (NIQE).
    \item We demonstrate that a model combining a Hybrid Swin Transformer with DINO and FFT losses (\texttt{DINO-Ranger}) achieves superior performance across most metrics, highlighting the power of self-supervised features.
    \item We uncover a crucial trade-off: while explicit degradation modeling improves robustness, it can slightly hinder the peak performance achievable with powerful perceptual losses, offering valuable guidance for future architectural design.
\end{itemize}

% =================================================================================================
% 3. RELATED WORK
% =================================================================================================
\section{Related Work}
\label{sec:related_work}

\subsection{Transformer-based SISR}
The success of Vision Transformers \cite{dosovitskiy2020image} in high-level vision tasks has inspired their application to image restoration. SwinIR \cite{liang2021swinir} adapted the Swin Transformer architecture, with its windowed and shifted-window self-attention, to achieve state-of-the-art results in SISR. The core idea is that local attention within windows captures texture details, while the shifting mechanism enables cross-window information flow, effectively modeling global context. Our work builds on this paradigm by implementing and comparing two powerful backbones: a \texttt{HybridSwinNet}, which fuses Swin blocks with convolutional layers, and a \texttt{TransformerESRGAN}, which integrates Residual Swin Transformer Blocks (RSTB) into a GAN framework.

\subsection{Degradation Modeling in SISR}
To address the poor generalization of models trained on idealized data, several works have proposed more realistic degradation pipelines. Real-ESRGAN \cite{wang2021real} introduced a "high-order" degradation process, sequentially applying blur, noise, resizing, and JPEG compression to create more challenging and realistic LR training pairs. This strategy forces the model to learn restoration in the presence of complex artifacts, rather than simply inverting a bicubic downsampling kernel. Our \texttt{NovelDegradationPipeline} follows this principle, incorporating a randomized sequence of corruptions, including chromatic aberration and sensor noise, to synthesize a diverse set of real-world-like degradations on-the-fly during training. Our \texttt{Degraded} and \texttt{DINO w/ Degraded} models directly leverage this pipeline.

\subsection{Perceptual and Self-Supervised Losses}
Pixel-wise losses like L1 or L2 often lead to overly smooth results, as they penalize perceptually plausible high-frequency details. Perceptual loss \cite{johnson2016perceptual}, which computes the distance between deep features extracted by a pre-trained network (typically VGG \cite{simonyan2014very}), was a major step forward. However, VGG features are inherently biased towards the classification task they were trained for.

Self-supervised learning (SSL) provides a powerful alternative. Models like DINO \cite{caron2021emerging} learn features by enforcing consistency between different augmented views of an image, without relying on any labels. These features have been shown to capture rich semantic and structural information. Recent work has started exploring SSL features for image restoration \cite{wang2023dall}. Our work is the first, to our knowledge, to systematically compare a DINO-based perceptual loss against other components in a comprehensive SISR study, demonstrating its significant advantages.

\begin{table*}[t!]
\centering
\caption{\textbf{Comparison of Architectural and Training Components Across the Four Proposed Models.} Each model represents a unique combination of a generator backbone, a training data generation strategy (degradation pipeline), and a loss function suite.}
\label{tab:model_components}
\resizebox{\textwidth}{!}{%
\begin{tabular}{@{}lcccc@{}}
\toprule
\textbf{Component} & \textbf{Degraded} & \textbf{DINO w/ Degraded} & \textbf{DINO-Ranger} & \textbf{Proposed-Merge} \\
\midrule
\textbf{Generator Backbone} & TransformerESRGAN & HybridSwinNet & HybridSwinNet & TransformerESRGAN \\
\textbf{Discriminator} & Multi-Scale & U-Net & U-Net & Multi-Scale \\
\midrule
\textbf{Degradation Pipeline} & \textbf{Yes} (\texttt{NovelDegradation}) & \textbf{Yes} (\texttt{NovelDegradation}) & No (Bicubic LR) & \textbf{Yes} (\texttt{NovelDegradation}) \\
\midrule
\textbf{Perceptual Loss} & VGG-based (Implicit) & \textbf{DINO} & \textbf{DINO} & \textbf{DINO} \\
\textbf{Frequency (FFT) Loss} & No & \textbf{Yes} & \textbf{Yes} & \textbf{Yes} \\
\textbf{Adversarial Loss} & Yes & Yes & Yes & Yes \\
\bottomrule
\end{tabular}%
}
\end{table*}

% =================================================================================================
% 4. PROPOSED METHODS
% =================================================================================================
\section{Proposed Methods}
\label{sec:methods}
We investigate the effectiveness of modern SISR components by constructing and comparing four distinct models. Each model is a unique amalgamation of a generator/discriminator architecture, a data degradation strategy, and a loss function suite. The key components are detailed below, and their specific combinations are summarized in Table \ref{tab:model_components}.

\subsection{Transformer-Based Generator Backbones}
Our study employs two powerful Transformer-based backbones.
\begin{itemize}
    \item \textbf{HybridSwinNet}: This generator uses a series of \texttt{HybridBlock} modules. Each block contains two Swin Transformer blocks—one with regular windowed attention (W-MSA) and one with shifted-window attention (SW-MSA)—followed by a residual convolutional layer. This design allows for both long-range dependency modeling via self-attention and effective local feature extraction via convolutions. It is used in our \texttt{DINO-Ranger} and \texttt{DINO w/ Degraded} models.
    \item \textbf{TransformerESRGAN}: This architecture, used in our \texttt{Degraded} and \texttt{Proposed-Merge} models, is built upon Residual Swin Transformer Blocks (RSTB). Each RSTB contains multiple Swin Transformer blocks and a final convolutional layer, with a residual connection around the entire block. This deep, residual structure allows for the stable training of very powerful restoration networks.
\end{itemize}

\subsection{Degradation-Aware Training}
To train models robust to real-world artifacts, we employ a \texttt{NovelDegradationPipeline} that generates LR images on-the-fly from HR source images. This pipeline applies a randomized sequence of the following degradations before final downsampling:
\begin{enumerate}
    \item \textbf{Chromatic Aberration}: Simulates lens imperfections by randomly shifting the R and B color channels.
    \item \textbf{Sensor Noise}: Adds a combination of Poisson (shot) and Gaussian (read) noise to mimic camera sensor characteristics.
    \item \textbf{Motion Blur}: Applies a random linear or rotational motion blur kernel to simulate camera shake.
    \item \textbf{JPEG Compression}: Compresses the image with a random quality factor (30-85) to introduce blocky artifacts.
\end{enumerate}
This pipeline ensures the network is exposed to a wide variety of realistic corruptions, preventing it from overfitting to a simple, clean downsampling kernel.

\subsection{Advanced Loss Formulation}
Our primary loss function is a weighted sum of four components: $L_{G} = w_{L1}L_{L1} + w_{p}L_{p} + w_{g}L_{g} + w_{fft}L_{fft}$.
\begin{itemize}
    \item \textbf{Pixel Loss ($L_{L1}$)}: The L1 distance between the super-resolved image $G(I_{LR})$ and the ground-truth HR image $I_{HR}$. It ensures pixel-level fidelity.
    \item \textbf{Adversarial Loss ($L_{g}$)}: A GAN loss that encourages the generator to produce images indistinguishable from real HR images. We use a multi-scale discriminator for our more advanced models.
    \item \textbf{DINO Perceptual Loss ($L_{p}$)}: Instead of a traditional VGG-based loss, we use features from a pre-trained, frozen DINO ViT-S/16 model. The loss is the L1 distance between the DINO features of the generated and ground-truth images:
    \begin{equation}
        L_{p} = ||\phi_{DINO}(G(I_{LR})) - \phi_{DINO}(I_{HR})||_1
    \end{equation}
    where $\phi_{DINO}$ is the feature extractor. DINO's self-supervised training allows it to capture high-level semantic and textural information without classification-specific biases.
    \item \textbf{Frequency (FFT) Loss ($L_{fft}$)}: This loss operates in the frequency domain. It minimizes the L1 distance between the magnitude of the Fast Fourier Transform (FFT) of the generated and ground-truth images. This encourages the model to restore high-frequency details (e.g., textures, fine edges) that are often lost with pixel-wise losses alone.
\end{itemize}

\subsection{Model Amalgamations}
We now define our four models based on the components above.
\begin{itemize}
    \item \textbf{Degraded}: Our baseline model. It uses the powerful \texttt{TransformerESRGAN} backbone and a multi-scale discriminator, trained with the \texttt{NovelDegradationPipeline}. Its loss function is a standard combination of L1, adversarial, and a VGG-based perceptual loss.
    \item \textbf{DINO w/ Degraded}: This model tests the direct injection of advanced losses into a degradation-aware framework. It uses the \texttt{HybridSwinNet} backbone, is trained with the \texttt{NovelDegradationPipeline}, and is optimized with our full loss suite, including DINO and FFT losses.
    \item \textbf{DINO-Ranger}: This model prioritizes architectural and loss function performance on cleaner data. It uses the \texttt{HybridSwinNet} backbone but is trained on LR images generated via simple bicubic downsampling. It is optimized with the full DINO and FFT loss suite.
    \item \textbf{Proposed-Merge}: This model represents the most comprehensive amalgamation. It combines the powerful \texttt{TransformerESRGAN} backbone and multi-scale discriminator with both the \texttt{NovelDegradationPipeline} and the full DINO/FFT loss suite, aiming for a model that is both robust and capable of high-fidelity reconstruction.
\end{itemize}

\begin{table*}[t!]
\centering
\caption{\textbf{Quantitative Results for x4 Super-Resolution on Standard Benchmarks.} We report PSNR (dB), SSIM, LPIPS, and NIQE. Higher PSNR/SSIM are better. Lower LPIPS/NIQE are better. Best results are in \textbf{bold}, second best are \underline{underlined}.}
\label{tab:quantitative_results}
\resizebox{\textwidth}{!}{%
\begin{tabular}{@{}l|cccc|cccc|cccc|cccc@{}}
\toprule
& \multicolumn{4}{c|}{\textbf{Set5}} & \multicolumn{4}{c|}{\textbf{Set14}} & \multicolumn{4}{c|}{\textbf{BSD100}} & \multicolumn{4}{c}{\textbf{Urban100}} \\
\cmidrule(l){2-5} \cmidrule(l){6-9} \cmidrule(l){10-13} \cmidrule(l){14-17}
\textbf{Model} & PSNR$\uparrow$ & SSIM$\uparrow$ & LPIPS$\downarrow$ & NIQE$\downarrow$ & PSNR$\uparrow$ & SSIM$\uparrow$ & LPIPS$\downarrow$ & NIQE$\downarrow$ & PSNR$\uparrow$ & SSIM$\uparrow$ & LPIPS$\downarrow$ & NIQE$\downarrow$ & PSNR$\uparrow$ & SSIM$\uparrow$ & LPIPS$\downarrow$ & NIQE$\downarrow$ \\
\midrule
Degraded & 28.36 & 0.8228 & 0.1144 & \underline{4.68} & 25.07 & 0.6683 & 0.1923 & \underline{3.91} & 24.81 & 0.6360 & 0.2089 & \textbf{3.53} & 22.95 & 0.6906 & 0.1774 & \underline{3.56} \\
DINO w/ Deg. & 29.09 & 0.8378 & 0.1150 & 5.10 & 26.38 & 0.7043 & 0.1775 & 3.97 & 25.60 & 0.6543 & 0.2072 & 3.66 & 24.14 & 0.7157 & 0.1801 & 3.64 \\
Proposed-Merge & \underline{30.25} & \underline{0.8651} & 0.1176 & 5.55 & \underline{26.78} & \underline{0.7234} & \underline{0.1724} & 4.00 & \underline{26.12} & \underline{0.6843} & \underline{0.2059} & \underline{3.50} & \underline{24.71} & \underline{0.7457} & \underline{0.1638} & 3.94 \\
DINO-Ranger & \textbf{30.95} & \textbf{0.8662} & \textbf{0.0877} & 5.18 & \textbf{27.36} & \textbf{0.7345} & \textbf{0.1458} & 3.98 & \textbf{26.32} & \textbf{0.6854} & \textbf{0.1890} & 3.37 & \textbf{25.00} & \textbf{0.7479} & \textbf{0.1519} & 3.71 \\
\bottomrule
\end{tabular}%
}
\end{table*}

% =================================================================================================
% 5. EXPERIMENTS
% =================================================================================================
\section{Experiments}
\label{sec:experiments}

\subsection{Experimental Setup}
\textbf{Datasets and Training}: All models were trained on the DF2K dataset, a combination of DIV2K \cite{agustsson2017ntire} and Flickr2K \cite{timofte2017ntire}. We evaluated performance on four standard benchmarks: Set5 \cite{bevilacqua2012low}, Set14 \cite{zeyde2010single}, BSD100 \cite{martin2001database}, and Urban100 \cite{huang2015single}. These datasets cover a range of content, from simple objects and faces (Set5) to diverse natural scenes (BSD100) and challenging structured patterns (Urban100). All evaluations were performed for x4 super-resolution.

\textbf{Evaluation Metrics}: We use four metrics to provide a comprehensive evaluation:
\begin{itemize}
    \item \textbf{PSNR} and \textbf{SSIM}: Widely used metrics to measure pixel-level fidelity and structural similarity. They are calculated on the Y channel of the YCbCr color space.
    \item \textbf{LPIPS} (Learned Perceptual Image Patch Similarity) \cite{zhang2018unreasonable}: A full-reference metric that measures perceptual similarity, correlating well with human judgment.
    \item \textbf{NIQE} (Natural Image Quality Evaluator) \cite{mittal2012making}: A no-reference metric that assesses the statistical "naturalness" of an image. Lower scores indicate the image statistics are closer to those of a pristine natural image.
\end{itemize}

\subsection{Quantitative Analysis}
Table \ref{tab:quantitative_results} presents the full quantitative results. The \texttt{DINO-Ranger} model consistently achieves the highest PSNR, SSIM, and LPIPS scores across all four datasets. This demonstrates the profound impact of combining a strong Transformer backbone with a superior perceptual loss. On the challenging Urban100 dataset, which is rich in geometric structures, \texttt{DINO-Ranger} outperforms the \texttt{Degraded} baseline by a remarkable 2.05 dB in PSNR.

The \texttt{Proposed-Merge} model, which integrates all advanced components, is a very strong runner-up, consistently securing the second-best scores in PSNR and SSIM. This confirms that its architecture, which fuses the powerful \texttt{TransformerESRGAN} with the full degradation and loss suite, is highly effective.

An interesting trend emerges with the NIQE metric. The \texttt{Degraded} model often achieves the best or second-best NIQE score, particularly on BSD100. This suggests that training explicitly on noisy, corrupted data guides the generator to produce outputs with more natural statistical properties, even if they are not the most faithful reconstructions. In contrast, the top-performing \texttt{DINO-Ranger} and \texttt{Proposed-Merge} models sometimes yield higher (worse) NIQE scores, a known phenomenon where GANs or powerful perceptual models can produce textures that are slightly "hyper-realistic" or "unnatural" to a statistical model.

The \texttt{DINO w/ Degraded} model consistently sits between the \texttt{Degraded} baseline and the top two performers. This result is crucial: simply adding DINO loss to a degradation-aware framework improves performance over the baseline, but it does not reach the heights of the \texttt{DINO-Ranger}. This implies a potential tension between the degradation pipeline, which introduces artifacts, and the DINO loss, which aims to match features of clean, pristine images.

%==============================================================================================================
% REVISED QUALITATIVE ANALYSIS SECTION
%==============================================================================================================
\subsection{Qualitative Analysis}
Figure \ref{fig:qualitative_grid} presents a comprehensive visual comparison on representative images from all four benchmark datasets, corroborating the quantitative findings in Table \ref{tab:quantitative_results}. The grid structure allows for a direct comparison of how each model performs on different types of image content.

The most striking visual improvement comes from the DINO-based models (\texttt{DINO-Ranger}, \texttt{Proposed-Merge}). On the Urban100 example (top row), they excel at restoring sharp, straight lines and fine geometric patterns on the building facade, which are heavily blurred in the LR input. The \texttt{Degraded} model, while producing a clean image, fails to recover these high-frequency structures, resulting in a softer appearance.

On the Set14 image containing text (second row), the advantage of the advanced loss suite is clear. Both \texttt{DINO-Ranger} and \texttt{Proposed-Merge} reconstruct significantly more legible characters compared to the other models, demonstrating their superior ability to restore fine, meaningful details.

For the natural scene from BSD100 (third row), the trade-off with degradation modeling becomes apparent. The \texttt{Degraded} model's output is visually pleasant and artifact-free, aligning with its excellent NIQE score. However, it smooths over the intricate fur texture. In contrast, \texttt{DINO-Ranger} restores a much sharper, more detailed texture that is closer to the ground truth, highlighting the power of the DINO loss for perceptual realism.

Finally, the result on the Set5 "baby" image (bottom row) shows that \texttt{DINO-Ranger} and \texttt{Proposed-Merge} successfully reconstruct subtle skin textures and hair strands, whereas the \texttt{Degraded} model produces a more uniform, slightly plastic-like skin texture. Overall, the visual evidence strongly supports the conclusion that the amalgamation of a Swin Transformer backbone with DINO and FFT losses provides the most perceptually convincing and high-fidelity results.

\begin{figure*}[t!]
    \centering
    % Define a common width for all subfigures
    \newcommand{\gridimgwidth}{0.155\textwidth}

    % --- Column Headers (First Row with Captions) ---
    \begin{subfigure}[b]{\gridimgwidth}
        \centering
        \includegraphics[width=\linewidth]{example-image-a} % Urban100 LR
        \caption*{\small LR Input (Urban100)}
    \end{subfigure}\hfill
    \begin{subfigure}[b]{\gridimgwidth}
        \centering
        \includegraphics[width=\linewidth]{example-image-b} % Urban100 Degraded
        \caption*{\small Degraded}
    \end{subfigure}\hfill
    \begin{subfigure}[b]{\gridimgwidth}
        \centering
        \includegraphics[width=\linewidth]{example-image-c} % Urban100 DINO w/ Deg
        \caption*{\small DINO w/ Degraded}
    \end{subfigure}\hfill
    \begin{subfigure}[b]{\gridimgwidth}
        \centering
        \includegraphics[width=\linewidth]{example-image} % Urban100 Proposed-Merge
        \caption*{\small Proposed-Merge}
    \end{subfigure}\hfill
    \begin{subfigure}[b]{\gridimgwidth}
        \centering
        \includegraphics[width=\linewidth]{example-image-a} % Urban100 DINO-Ranger
        \caption*{\small DINO-Ranger}
    \end{subfigure}\hfill
    \begin{subfigure}[b]{\gridimgwidth}
        \centering
        \includegraphics[width=\linewidth]{example-image-b} % Urban100 GT
        \caption*{\small Ground Truth}
    \end{subfigure}

    \vspace{1.5mm} % Space between rows

    % --- Second Row (Set14) ---
    \begin{subfigure}[b]{\gridimgwidth} \includegraphics[width=\linewidth]{example-image-c} \caption*{\small LR Input (Set14)} \end{subfigure}\hfill
    \begin{subfigure}[b]{\gridimgwidth} \includegraphics[width=\linewidth]{example-image} \end{subfigure}\hfill
    \begin{subfigure}[b]{\gridimgwidth} \includegraphics[width=\linewidth]{example-image-a} \end{subfigure}\hfill
    \begin{subfigure}[b]{\gridimgwidth} \includegraphics[width=\linewidth]{example-image-b} \end{subfigure}\hfill
    \begin{subfigure}[b]{\gridimgwidth} \includegraphics[width=\linewidth]{example-image-c} \end{subfigure}\hfill
    \begin{subfigure}[b]{\gridimgwidth} \includegraphics[width=\linewidth]{example-image} \end{subfigure}
    
    \vspace{1.5mm} % Space between rows

    % --- Third Row (BSD100) ---
    \begin{subfigure}[b]{\gridimgwidth} \includegraphics[width=\linewidth]{example-image-a} \caption*{\small LR Input (BSD100)} \end{subfigure}\hfill
    \begin{subfigure}[b]{\gridimgwidth} \includegraphics[width=\linewidth]{example-image-b} \end{subfigure}\hfill
    \begin{subfigure}[b]{\gridimgwidth} \includegraphics[width=\linewidth]{example-image-c} \end{subfigure}\hfill
    \begin{subfigure}[b]{\gridimgwidth} \includegraphics[width=\linewidth]{example-image} \end{subfigure}\hfill
    \begin{subfigure}[b]{\gridimgwidth} \includegraphics[width=\linewidth]{example-image-a} \end{subfigure}\hfill
    \begin{subfigure}[b]{\gridimgwidth} \includegraphics[width=\linewidth]{example-image-b} \end{subfigure}
    
    \vspace{1.5mm} % Space between rows
    
    % --- Fourth Row (Set5) ---
    \begin{subfigure}[b]{\gridimgwidth} \includegraphics[width=\linewidth]{example-image-c} \caption*{\small LR Input (Set5)} \end{subfigure}\hfill
    \begin{subfigure}[b]{\gridimgwidth} \includegraphics[width=\linewidth]{example-image} \end{subfigure}\hfill
    \begin{subfigure}[b]{\gridimgwidth} \includegraphics[width=\linewidth]{example-image-a} \end{subfigure}\hfill
    \begin{subfigure}[b]{\gridimgwidth} \includegraphics[width=\linewidth]{example-image-b} \end{subfigure}\hfill
    \begin{subfigure}[b]{\gridimgwidth} \includegraphics[width=\linewidth]{example-image-c} \end{subfigure}\hfill
    \begin{subfigure}[b]{\gridimgwidth} \includegraphics[width=\linewidth]{example-image} \end{subfigure}

    \caption{\textbf{Qualitative comparison on representative images from all four benchmark datasets for x4 super-resolution.} The columns show the bicubic LR input, the outputs from our four proposed models, and the ground truth. Each row corresponds to a different dataset. The DINO-based models, especially \texttt{DINO-Ranger} and \texttt{Proposed-Merge}, consistently restore sharper textures and more accurate geometric structures (e.g., Urban100 building facade, Set14 text). The \texttt{Degraded} model produces visually natural but softer results. \textit{(Note: Placeholder images are used. Replace with actual output images from your experiments.)}}
    \label{fig:qualitative_grid}
\end{figure*}

% =================================================================================================
% 6. DISCUSSION
% =================================================================================================
\section{Discussion}
Our comparative study reveals several key insights into the design of modern SISR systems.

\textbf{The Dominance of Self-Supervised Perceptual Loss.} The consistently superior performance of the models using DINO loss (\texttt{DINO-Ranger}, \texttt{Proposed-Merge}) in terms of PSNR, SSIM, and LPIPS is the most significant finding. It strongly suggests that features learned via self-supervision are more effective for guiding image restoration than features from classification-oriented networks like VGG. The DINO features appear to capture a more nuanced understanding of texture and structure, enabling the generator to produce reconstructions that are both quantitatively accurate and perceptually pleasing.

\textbf{The Degradation Modeling Trade-Off.} While training with a realistic degradation pipeline is crucial for robustness, our results highlight a subtle trade-off. The \texttt{Degraded} model achieves excellent NIQE scores, indicating high statistical naturalness. However, the top-performing models, particularly \texttt{DINO-Ranger}, which was trained on cleaner bicubic data, achieved higher fidelity. This suggests that when the goal is to maximize reconstruction quality for a known (or simple) degradation, training on heavily corrupted data might be a form of over-regularization that prevents the model from achieving peak performance. The ideal strategy may involve curriculum learning, where the severity of degradation is adjusted during training.

\textbf{Architectural Synergy.} Both the \texttt{HybridSwinNet} and \texttt{TransformerESRGAN} backbones proved to be highly capable. The success of the \texttt{Proposed-Merge} model shows that combining a powerful architecture with a sophisticated degradation pipeline and an advanced loss suite can create an extremely competitive "all-in-one" solution. The performance difference between \texttt{Proposed-Merge} and \texttt{DINO-Ranger} is small, indicating that both backbone architectures are near the state of the art, and the choice between them may depend on specific implementation or efficiency considerations.

% =================================================================================================
% 7. CONCLUSION
% =================================================================================================
\section{Conclusion}
\label{sec:conclusion}
In this work, we presented a comprehensive comparative study of four SISR models, each an amalgamation of modern architectural and training components. By systematically evaluating the roles of Transformer backbones, realistic degradation modeling, and an advanced loss suite featuring self-supervised DINO features, we have provided clear insights into their individual and synergistic effects. Our results show that a model combining a Hybrid Swin Transformer with DINO and FFT losses achieves state-of-the-art results in x4 super-resolution, significantly outperforming baselines in both quantitative and perceptual metrics. We also identified a key design trade-off between maximizing reconstruction fidelity and optimizing for statistical naturalness via degradation modeling. This study underscores the immense potential of self-supervised features in image restoration and offers a data-driven guide for practitioners aiming to build the next generation of robust and powerful SISR systems.

% =================================================================================================
% REFERENCES
% =================================================================================================
{\small
\bibliographystyle{ieee_fullname}
% \bibliography{egbib}
}

\begin{thebibliography}{10}\itemsep=-1pt

\bibitem{dong2015image}
Chao Dong, Chen~Change Loy, Kaiming He, and Xiaoou Tang.
\newblock Image super-resolution using deep convolutional networks.
\newblock {\em IEEE transactions on pattern analysis and machine intelligence}, 38(2):295--307, 2015.

\bibitem{lim2017enhanced}
Bee Lim, Sanghyun Son, Heewon Kim, Seungjun Nah, and Kyoung~Mu Lee.
\newblock Enhanced deep residual networks for single image super-resolution.
\newblock In {\em Proceedings of the IEEE conference on computer vision and pattern recognition workshops}, pages 136--144, 2017.

\bibitem{dosovitskiy2020image}
Alexey Dosovitskiy, Lucas Beyer, Alexander Kolesnikov, Dirk Weissenborn, Xiaohua Zhai, Thomas Unterthiner, Mostafa Dehghani, Matthias Minderer, Georg Heigold, Sylvain Gelly, et~al.
\newblock An image is worth 16x16 words: Transformers for image recognition at scale.
\newblock {\em arXiv preprint arXiv:2010.11929}, 2020.

\bibitem{liang2021swinir}
Jingyun Liang, Jiezhang Cao, Guolei Sun, Kai Zhang, Luc Van~Gool, and Radu Timofte.
\newblock Swinir: Image restoration using swin transformer.
\newblock In {\em Proceedings of the IEEE/CVF international conference on computer vision}, pages 1833--1844, 2021.

\bibitem{wang2021real}
Xintao Wang, Liangbin Xie, Chao Dong, and Ying Shan.
\newblock Real-esrgan: Training real-world blind super-resolution with pure synthetic data.
\newblock In {\em Proceedings of the IEEE/CVF international conference on computer vision}, pages 1905--1914, 2021.

\bibitem{simonyan2014very}
Karen Simonyan and Andrew Zisserman.
\newblock Very deep convolutional networks for large-scale image recognition.
\newblock {\em arXiv preprint arXiv:1409.1556}, 2014.

\bibitem{caron2021emerging}
Mathilde Caron, Hugo Touvron, Ishan Misra, Herv{\'e} J{\'e}gou, Julien Mairal, Piotr Bojanowski, and Armand Joulin.
\newblock Emerging properties in self-supervised vision transformers.
\newblock In {\em Proceedings of the IEEE/CVF international conference on computer vision}, pages 9650--9660, 2021.

\bibitem{johnson2016perceptual}
Justin Johnson, Alexandre Alahi, and Li Fei-Fei.
\newblock Perceptual losses for real-time style transfer and super-resolution.
\newblock In {\em European conference on computer vision}, pages 694--711. Springer, 2016.

\bibitem{wang2023dall}
Guanhua Wang, Jian-wei Zhang, and Zhen-zhong Chen.
\newblock Dall-sr: A diffusion-based super-resolution model with a lightweight conditional adapter.
\newblock {\em arXiv preprint arXiv:2311.16904}, 2023.

\bibitem{agustsson2017ntire}
Eirikur Agustsson and Radu Timofte.
\newblock Ntire 2017 challenge on single image super-resolution: Dataset and study.
\newblock In {\em Proceedings of the IEEE conference on computer vision and pattern recognition workshops}, pages 126--135, 2017.

\bibitem{timofte2017ntire}
Radu Timofte, Eirikur Agustsson, Luc Van Gool, Ming-Hsuan Yang, Lei Zhang, et~al.
\newblock Ntire 2017 challenge on single image super-resolution: Methods and results.
\newblock In {\em 2017 IEEE conference on computer vision and pattern recognition workshops (CVPRW)}, pages 1110--1121. IEEE, 2017.

\bibitem{bevilacqua2012low}
Marco Bevilacqua, Aline Roumy, Christine Guillemot, and Marie-Line Alberi-Morel.
\newblock Low-complexity single-image super-resolution based on nonnegative neighbor embedding.
\newblock In {\em British Machine Vision Conference}, 2012.

\bibitem{zeyde2010single}
Roman Zeyde, Michael Elad, and Matan Protter.
\newblock On single image scale-up using sparse-representations.
\newblock In {\em International conference on curves and surfaces}, pages 711--730. Springer, 2010.

\bibitem{martin2001database}
David Martin, Charless Fowlkes, Doron Tal, and Jitendra Malik.
\newblock A database of human segmented natural images and its application to evaluating segmentation algorithms and measuring ecological statistics.
\newblock In {\em Proceedings of the eighth ieee international conference on computer vision}, volume~2, pages 416--423. IEEE, 2001.

\bibitem{huang2015single}
Jia-Bin Huang, Abhishek Singh, and Narendra Ahuja.
\newblock Single image super-resolution from transformed self-exemplars.
\newblock In {\em Proceedings of the IEEE conference on computer vision and pattern recognition}, pages 5197--5206, 2015.

\bibitem{zhang2018unreasonable}
Richard Zhang, Phillip Isola, Alexei~A Efros, Eli Shechtman, and Oliver Wang.
\newblock The unreasonable effectiveness of deep features as a perceptual metric.
\newblock In {\em Proceedings of the IEEE conference on computer vision and pattern recognition}, pages 586--595, 2018.

\bibitem{mittal2012making}
Anish Mittal, Anush~Krishna Moorthy, and Alan~Conrad Bovik.
\newblock Making a “completely blind” image quality analyzer.
\newblock {\em IEEE Signal processing letters}, 20(3):209--212, 2012.

\end{thebibliography}

\end{document}